\documentclass[11pt,a4paper]{article}
% Estilo comum dos PDFs individuais da Sprint 1.
\usepackage[utf8]{inputenc}
\usepackage[T1]{fontenc}
\usepackage[brazilian]{babel}
\usepackage[a4paper,margin=2.35cm]{geometry}
\usepackage{xcolor}
\usepackage{hyperref}
\usepackage{enumitem}
\usepackage{listings}
\usepackage{parskip}
\usepackage{needspace}
\definecolor{azul}{RGB}{20,63,112}
\hypersetup{colorlinks=true,linkcolor=azul,urlcolor=azul}
\setlist[itemize]{leftmargin=*,itemsep=3pt}
\lstset{basicstyle=\ttfamily\small,frame=single,breaklines=true,columns=fullflexible}
\sloppy
\newcommand{\codigo}[1]{\texttt{#1}}
\newcommand{\titulo}[2]{\begin{center}{\LARGE\bfseries #1\par}\vspace{4pt}{\large #2\par}\vspace{8pt}\rule{0.8\textwidth}{0.4pt}\par\vspace{5pt}{\small Matheus Moreli Durão Silva -- Matrícula 12621CBS217\par Cibersegurança - Bacharelado - Noturno\par}\end{center}}
\newcommand{\cardfonte}[1]{\par\smallskip\noindent\textit{Referência: }\href{#1}{Trello - enunciado e checklist do card}.}

\begin{document}\titulo{S1\_15 - Tratamento de erros}{Salvar e carregar ativos com segurança}
\section*{Entregável correto}
O exercício solicitado é um sistema de cadastro de ativos com tratamento de erros na leitura e escrita de arquivos. Ele também pede validar entrada numérica e usar \codigo{raise} para dado inválido. Erro não é algo a esconder: é uma situação prevista que precisa receber uma resposta compreensível.
\section*{Como \codigo{try/except} funciona}
Dentro de \codigo{try} fica a operação que pode falhar. \codigo{except ValueError} trata, por exemplo, texto onde era esperado inteiro; \codigo{except FileNotFoundError} trata arquivo ausente. \codigo{else} roda apenas se não houver falha; \codigo{finally} roda sempre, sendo útil para encerramentos. Tratar exceções específicas é melhor que usar \codigo{except:} sem nome, que pode esconder problemas reais.
\begin{lstlisting}
class DadosInvalidosError(Exception):
    pass

def validar_id(texto):
    if not texto.isdigit():
        raise DadosInvalidosError("ID deve conter somente numeros.")
    return int(texto)

try:
    id_ativo = validar_id(input("ID: "))
except DadosInvalidosError as erro:
    print(f"Erro: {erro}")
\end{lstlisting}
\section*{Aplicação no projeto}
No \codigo{app.py}, a validação impede cadastro inconsistente; a leitura do JSON distingue arquivo inexistente de conteúdo inválido; o salvamento informa se conseguiu concluir. Uma exceção personalizada comunica melhor a regra quebrada do que uma mensagem genérica.
\section*{Checklist demonstrado}
Erro versus exceção, \codigo{try/except}, \codigo{else/finally}, exceções específicas, \codigo{raise}, validação e erros de operação foram estudados. Testes: digitar ID com letras, tentar carregar arquivo inexistente e salvar dados válidos.
\cardfonte{https://trello.com/c/qSZwgosj}
\end{document}
